\begin{figure}[htbp]
\centering
\begin{tikzpicture}[
    font=\footnotesize,
    >=Stealth,
    instrument/.style={
        rectangle, rounded corners=4pt, draw=purple!70!black,
        fill=purple!7, minimum width=31mm, minimum height=12mm,
        align=center, font=\bfseries
    },
    hypothesis/.style={
        rectangle, rounded corners=3pt, minimum width=34mm,
        minimum height=11mm, align=center, font=\scriptsize
    },
    evidence/.style={
        rectangle, rounded corners=5pt, draw=orange!80!black,
        fill=orange!12, minimum width=37mm, minimum height=15mm,
        text width=37mm, align=center, font=\bfseries
    },
    posterior/.style={
        rectangle, rounded corners=3pt, draw=blue!65!black,
        fill=blue!6, minimum width=43mm, minimum height=52mm,
        align=left
    },
    flow/.style={draw=black!55, thick, -{Stealth[length=2.2mm]}},
    likelihood/.style={
        fill=white, inner sep=1.2pt, font=\tiny\bfseries
    },
    note/.style={align=center, font=\scriptsize}
]

% Instrumento y roca observada
\node[instrument] (nir) at (-5.05,2.95)
    {Espectrómetro NIR\\del rover};

\draw[black!65, thick] (-5.75,2.18) -- (-5.35,2.40) -- (-4.65,2.40);
\draw[fill=gray!25, draw=black!60, thick]
    (-3.95,2.16) -- (-3.55,2.48) -- (-3.05,2.32)
    -- (-2.82,1.88) -- (-3.25,1.67) -- (-3.78,1.77) -- cycle;
\node[font=\scriptsize, text=black!70] at (-3.38,1.47) {afloramiento rocoso};

\draw[orange!85!black, line width=1.1pt, -{Stealth[length=2mm]}]
    (nir.south east) -- (-3.66,2.26);
\draw[orange!55, dashed, thick]
    (-3.55,2.22) -- (-2.30,2.68);

% Espectro con absorción a 1 micrómetro
\begin{scope}[shift={(-1.95,1.95)},x=1.05cm,y=0.82cm]
    \fill[blue!3, rounded corners=3pt] (-0.25,-0.30) rectangle (4.15,2.25);
    \draw[black!45,->] (0,0) -- (3.85,0)
        node[below, font=\scriptsize] {$\lambda\;(\mu\mathrm{m})$};
    \draw[black!45,->] (0,0) -- (0,2.0)
        node[above, rotate=90, anchor=south, font=\scriptsize] {reflectancia};

    \fill[blue!16]
        plot[smooth] coordinates {
            (0.10,1.52) (0.55,1.60) (1.00,1.48)
            (1.45,1.20) (1.75,0.34) (2.05,1.15)
            (2.55,1.43) (3.10,1.53) (3.60,1.49)
        } -- (3.60,0) -- (0.10,0) -- cycle;
    \draw[blue!75!black, very thick]
        plot[smooth] coordinates {
            (0.10,1.52) (0.55,1.60) (1.00,1.48)
            (1.45,1.20) (1.75,0.34) (2.05,1.15)
            (2.55,1.43) (3.10,1.53) (3.60,1.49)
        };
    \draw[red!75!black, densely dashed, thick] (1.75,0) -- (1.75,1.82);
    \node[font=\tiny] at (0.62,-0.20) {$0.7$};
    \node[font=\tiny, text=red!75!black] at (1.75,-0.20) {$1.0$};
    \node[font=\tiny] at (2.90,-0.20) {$1.3$};
    \node[font=\scriptsize\bfseries, text=red!75!black, align=center]
        at (2.63,0.48) {banda de\\absorción};
\end{scope}

\node[note, text=purple!70!black] at (0.12,4.05)
    {Medición observacional};

% Evidencia central
\node[evidence] (band) at (0,-1.05)
    {$B$: banda profunda en $1\,\mu\mathrm{m}$\\firma asociada al piroxeno};
\draw[flow] (0,1.72) -- (band.north);

% Hipótesis y probabilidades a priori
\node[hypothesis, draw=blue!70!black, fill=blue!9] (a1) at (-4.75,0.15)
    {$A_1$: roca basáltica\\$P(A_1)=0.20$};
\node[hypothesis, draw=green!50!black, fill=green!9] (a2) at (-4.75,-1.25)
    {$A_2$: condrita carbonácea\\$P(A_2)=0.60$};
\node[hypothesis, draw=gray!70!black, fill=gray!12] (a3) at (-4.75,-2.65)
    {$A_3$: origen metálico\\$P(A_3)=0.20$};

\node[font=\scriptsize\bfseries, text=black!65] at (-4.75,0.92)
    {HIPÓTESIS A PRIORI};

\draw[flow, blue!70!black] (a1.east) -- ++(0.55,0) |- ([yshift=5mm]band.west);
\draw[flow, green!45!black] (a2.east) -- (band.west);
\draw[flow, gray!70!black] (a3.east) -- ++(0.55,0) |- ([yshift=-5mm]band.west);
\node[likelihood, text=blue!70!black] at (-2.50,0.30) {$P(B\mid A_1)=0.90$};
\node[likelihood, text=green!45!black] at (-2.50,-1.02) {$P(B\mid A_2)=0.10$};
\node[likelihood, text=gray!70!black] at (-2.50,-2.35) {$P(B\mid A_3)=0.05$};

% Resultado posterior
\node[posterior] (post) at (4.55,-1.15) {};
\node[font=\scriptsize\bfseries, text=blue!65!black]
    at (4.55,1.18) {POSTERIOR DADO $B$};

\draw[flow] (band.east) -- (post.west)
    node[likelihood, above] {Bayes};

\node[anchor=west, font=\scriptsize] at (2.75,0.40) {$P(A_1\mid B)=0.72$};
\fill[blue!65] (2.75,0.10) rectangle ++(2.70,0.22);

\node[anchor=west, font=\scriptsize] at (2.75,-0.55) {$P(A_2\mid B)=0.24$};
\fill[green!55!black] (2.75,-0.85) rectangle ++(0.90,0.22);

\node[anchor=west, font=\scriptsize] at (2.75,-1.50) {$P(A_3\mid B)=0.04$};
\fill[gray!65] (2.75,-1.80) rectangle ++(0.15,0.22);

\node[
    rectangle, rounded corners=3pt, draw=blue!75!black,
    fill=blue!12, font=\scriptsize\bfseries, align=center,
    inner sep=3pt
] at (4.55,-2.55)
    {Clasificación MAP:\\$A_1$ (roca basáltica)};

% Normalización y lectura física
\node[
    note, draw=orange!60!black, fill=orange!5,
    rounded corners=3pt, inner sep=3pt, text width=43mm
] at (0,-3.55)
    {$P(B)=0.90(0.20)+0.10(0.60)+0.05(0.20)$\\
     $P(B)=0.25$};

\end{tikzpicture}
\caption{Esquema de clasificación mineralógica mediante espectroscopía de reflexión. El espectrómetro detecta una banda de absorción en $1\,\mu\mathrm{m}$ y combina las probabilidades a priori con las verosimilitudes mediante el Teorema de Bayes. La evidencia favorece a la hipótesis basáltica $A_1$, cuya probabilidad posterior es $0.72$.}
\end{figure}